\section{Conclusion}

This study investigated how to improve the reliability of LLM \emph{agent} tool calling
through lightweight, model-agnostic interventions rather than model retraining. We introduced three
mechanisms: tool-dependency skills, gated skills, and
actor-critic verification and evaluated them through a controlled ablation across WorkBench and API-Bank benchmarks
 and on three models. On WorkBench, averaged over the six domains, the integrated approach raised
GPT-4o-mini accuracy from 55.1\% to 64.0\% (a 16.2\% relative improvement) while reducing
harmful side effects. Crucially, every mechanism operates purely at the prompt and orchestration
level, making it a generic layer that can be attached to any tool-calling agent.

The approach is also inexpensive. Because the mechanisms add only a small amount of context and,
at most, a single auxiliary classifier or critic call per query, the total cost remained on the
order of one to two cents per task on GPT-4o-mini and a fraction of a cent on the open
7--8B models, a negligible overhead relative to the reliability gains, and orders of magnitude
cheaper than retraining model. This favorable cost profile, combined with
the training-free design, makes the method practical to deploy in real settings.

TODO: update numbers. Among the individual mechanisms, gated skills proved the most effective single intervention,
raising accuracy on every non-saturated domain for example, from 65.6\% to 84.4\% on email
and from 52.5\% to 67.5\% on project management with GPT-4o-mini by injecting only the
guidance relevant to each query rather than a fixed instruction set. Actor-critic verification
contributed primarily to safety, roughly halving the unwanted side-effect rate (e.g., 25.0\% to
12.5\% on project management) while adding modest accuracy, whereas tool-dependency skills helped
selectively, chiefly on tasks with genuine producer-consumer chains. Combining all three
mechanisms was not strictly additive: the full configuration excelled only in the hardest,
mixed-domain settings (e.g., $+11.6$ on analytics, $+4.3$ on multi-domain) and elsewhere trailed
gated skills alone. Improvements scaled with available headroom, vanishing on already-saturated
domains such as calendar.

Several limitations remain. Side-effect errors persist and. Gains varied
substantially across domains and model scales, suggesting that model and domain-specific tuning
may be required for maximal reliability. Future work should explore reinforcement learning and
process reward models targeted specifically at minimizing side effects while preserving
tool-calling reliability, dynamic error-recovery mechanisms that detect and reverse unintended
actions during execution, scalability to larger tool repertoires, and the security and
data-integrity risks that accompany increasingly autonomous tool use. 

In summary, this work shows that a generic, training-free, and inexpensive layer of tool-use interventions,
 a significant step toward developing more reliable and 
capable LLM agents for real-world applications by narrowing the gap between the 
theoretical potential of tool-augmented LLMs and their practical ability to execute complex tasks.
